% Source PDF pages 55--60; manual pages 2-26--2-31.
\subsection{Wind Cartesian Coordinate System}\label{sec:wind-coordinate}

Aerodynamic forces are usually functions of aerodynamic angles such as angle of
attack and sideslip. The wind coordinate system relates these aerodynamic angles to
the body-fixed system. Its origin is at the vehicle's center of mass, and its $x$ axis
is aligned with the air-relative velocity-system axis $\hat x_{V_w}$ and with
$\vec V_{w\oplus}$.

The wind-system $y$ and $z$ axes are obtained by rotating the corresponding
velocity-system axes through the geocentric bank angle $\mu_{gc}$ or geodetic bank
angle $\mu_{gd}$, as shown in \cref{fig:bank-angle-wind-system}. The two bank
angles differ slightly because the geocentric and geodetic local-horizon systems are
not identical, but both constructions produce the same wind system.

\begin{figure}[htbp]
  \centering
  \resizebox{0.76\linewidth}{!}{\begin{tikzpicture}[>=Latex,line cap=round,line join=round,font=\sffamily\scriptsize]
  \coordinate (O) at (0,0);

  % Reference local-horizon and velocity axes.
  \draw[->,line width=1.0pt] (O)--(3.15,0) node[right] {$\hat{y}_{V}$};
  \draw[->,line width=1.0pt] (O)--(0,-2.05) node[below] {$\hat{z}_{V}$};
  \draw[->,line width=1.0pt] (O)--(-0.72,-1.55) node[below left] {$\hat{z}_{w}$};
  \draw[->,line width=1.05pt] (O)--(2.45,-0.82) node[right] {$\hat{y}_{w}$};
  \draw[->,line width=0.9pt] (O)--(0.20,1.55) node[above] {$\hat{x}_{V}=\hat{x}_{w}$};

  % Bank angle arc.
  \draw[->] (1.35,0) arc[start angle=0,end angle=-18.5,radius=1.35];
  \node at (1.65,-0.25) {$\mu_{gc}$ or $\mu_{gd}$};

  % Measurement-plane annotation, moved clear of all vector labels.
  \node[align=left,anchor=east] at (-1.05,1.25)
    {Plane of measurement\\for sideslip angle $\beta_E$};
  \draw[->,line width=0.75pt] (-0.92,0.95)--(-0.08,0.35);

  \node[align=right,anchor=east] at (-0.55,0.05)
    {Geocentric or geodetic\\local horizontal};
  \draw[line width=0.75pt] (-0.42,0)--(O);
  \node[align=left,anchor=west] at (0.55,1.18)
    {$\hat{y}_{V}$ and $\hat{z}_{V}$ lie in the page;\\$\hat{x}_{V}$ is normal to the page.};
\end{tikzpicture}
}
  \caption{Bank angle and wind coordinate system.}
  \label{fig:bank-angle-wind-system}
\end{figure}

The wind unit vectors are
\begin{equation}
  \begin{aligned}
    \hat x_w &= \hat x_V,\\
    \hat y_w &= \cos\mu\,\hat y_V+\sin\mu\,\hat z_V,\\
    \hat z_w &= -\sin\mu\,\hat y_V+\cos\mu\,\hat z_V,
  \end{aligned}
  \label{eq:wind-unit-vectors}
\end{equation}
and are evaluated by \texttt{wind\_unit\_vectors}. The bank angle $\mu$ must be
referenced to the same system, geocentric or geodetic, as the velocity-system unit
vectors.

The wind and body systems are related through one of three equivalent pairs of
aerodynamic angles:
\begin{itemize}[leftmargin=1.6em]
  \item angle of attack $\alpha$ and Euler sideslip angle $\beta_E$;
  \item angle of attack $\alpha$ and sideslip angle $\beta$; or
  \item total angle of attack $\alpha_T$ and windward meridian $\phi_w$.
\end{itemize}

\taossubhead{Angle of attack and Euler sideslip angle}

As illustrated in \cref{fig:alpha-betae-definitions}, a rotation about the body
$-\hat y_b$ axis through the angle of attack $\alpha$, followed by a rotation about
the wind $\hat z_w$ axis through the Euler sideslip angle $\beta_E$, aligns the
body and wind axes. This rotational definition is commonly used for aircraft.
Angle of attack generally lies between $-180^\circ$ and $+180^\circ$, while Euler
sideslip generally lies between $-90^\circ$ and $+90^\circ$; because both are
rotation angles, however, the mathematics permits any values.

\begin{figure}[htbp]
  \centering
  \resizebox{0.78\linewidth}{!}{\begin{tikzpicture}[scale=0.76,>=Latex,line cap=round,line join=round]
  \coordinate (A) at (0,0); \coordinate (B) at (5.2,0);
  \coordinate (C) at (6.5,0.9); \coordinate (D) at (1.3,0.9);
  \coordinate (At) at (0,2.2); \coordinate (Bt) at (5.2,2.2);
  \coordinate (Ct) at (6.5,3.1); \coordinate (Dt) at (1.3,3.1);
  \draw (A)--(B)--(C)--(D)--cycle;
  \draw (At)--(Bt)--(Ct)--(Dt)--cycle;
  \draw (A)--(At) (B)--(Bt) (C)--(Ct) (D)--(Dt);
  \draw[dashed] (D)--(B) (D)--($(B)!0.50!(Bt)$);
  \draw[->,line width=1.1pt] (Dt)--(B) node[pos=0.85,below] {$\vec V_{w\oplus}$};
  \draw[->] (C)--($(C)+(1.35,0)$) node[right] {$\hat{x}_b$};
  \draw[->] (D)--($(D)+(-1.25,-0.75)$) node[left] {$\hat{y}_b$};
  \draw[->] (D)--($(D)+(0,-1.55)$) node[below] {$\hat{z}_b$};
  \draw[->] (B)--($(B)+(1.30,-0.90)$) node[right] {$\hat{x}_w$};
  \node at (4.15,2.70) {$+\alpha$};
  \node at (5.55,2.55) {$+\beta$};
  \node at (4.45,1.45) {$+\beta_E$};
\end{tikzpicture}
}
  \caption{Angle-of-attack and sideslip definitions for forward flight.}
  \label{fig:alpha-betae-definitions}
\end{figure}

The body unit vectors are related to the wind unit vectors by
\begin{equation}
  \begin{bmatrix}
    \hat x_b\\[0.2em]
    \hat y_b\\[0.2em]
    \hat z_b
  \end{bmatrix}
  =
  \TaosPitchMatrix{\alpha}
  \TaosEulerSideslipMatrix{\beta_E}
  \begin{bmatrix}
    \hat x_w\\[0.2em]
    \hat y_w\\[0.2em]
    \hat z_w
  \end{bmatrix}.
  \label{eq:wind-to-body-aero-transformation}
\end{equation}
Expanding the matrices gives
\begin{equation}
  \hat x_b
  =\cos\alpha\cos\beta_E\,\hat x_w
  -\cos\alpha\sin\beta_E\,\hat y_w
  -\sin\alpha\,\hat z_w,
  \label{eq:body-x-from-aero-angles}
\end{equation}
\begin{equation}
  \hat y_b
  =\sin\beta_E\,\hat x_w
  +\cos\beta_E\,\hat y_w,
  \label{eq:body-y-from-aero-angles}
\end{equation}
\begin{equation}
  \hat z_b
  =\sin\alpha\cos\beta_E\,\hat x_w
  -\sin\alpha\sin\beta_E\,\hat y_w
  +\cos\alpha\,\hat z_w.
  \label{eq:body-z-from-aero-angles}
\end{equation}

\taossubhead{Angle of attack and sideslip angle}

The non-Euler sideslip angle $\beta$ is the angle between $\hat x_b$ and the
projection of $\vec V_{w\oplus}$ onto the $\hat x_b\hat y_b$ plane. With this
sideslip definition, $\alpha$ is the angle between $\hat x_b$ and the projection of
$\vec V_{w\oplus}$ onto the $\hat x_b\hat z_b$ plane:
\begin{equation}
  \tan\alpha
  =\frac{\hat x_w\mathbin{\cdot}\hat z_b}
  {\hat x_w\mathbin{\cdot}\hat x_b},
  \label{eq:projected-angle-of-attack}
\end{equation}
\begin{equation}
  \tan\beta
  =\frac{\hat x_w\mathbin{\cdot}\hat y_b}
  {\hat x_w\mathbin{\cdot}\hat x_b}.
  \label{eq:projected-sideslip}
\end{equation}
These projection-based definitions are common for axisymmetric vehicles such as
reentry bodies and artillery shells. They are associated naturally with total angle of
attack and windward meridian. Because the angles are defined by projections rather
than independent rotations, not every combination of $\alpha$ and $\beta$ is valid.

When $\hat x_w\mathbin{\cdot}\hat x_b>0$, the vehicle is moving in the forward
direction and both $|\alpha|$ and $|\beta|$ are less than $90^\circ$, as in
\cref{fig:alpha-betae-definitions}. When
$\hat x_w\mathbin{\cdot}\hat x_b<0$, the vehicle is moving backward and both
angles have magnitudes between $90^\circ$ and $180^\circ$, as shown in
\cref{fig:alpha-beta-backward}.

\begin{figure}[htbp]
  \centering
  \resizebox{0.82\linewidth}{!}{\begin{tikzpicture}[scale=0.76,>=Latex,line cap=round,line join=round]
  \coordinate (A) at (0,0); \coordinate (B) at (5.2,0);
  \coordinate (C) at (6.5,0.9); \coordinate (D) at (1.3,0.9);
  \coordinate (At) at (0,2.2); \coordinate (Bt) at (5.2,2.2);
  \coordinate (Ct) at (6.5,3.1); \coordinate (Dt) at (1.3,3.1);
  \draw (A)--(B)--(C)--(D)--cycle;
  \draw (At)--(Bt)--(Ct)--(Dt)--cycle;
  \draw (A)--(At) (B)--(Bt) (C)--(Ct) (D)--(Dt);
  \draw[dashed] (D)--(B) (D)--($(B)!0.50!(Bt)$);
  \draw[->,line width=1.1pt] (Ct)--(A) node[pos=0.75,below] {$\vec V_{w\oplus}$};
  \draw[->] (Ct)--($(Ct)+(1.35,0)$) node[right] {$\hat{x}_b$};
  \draw[->] (B)--($(B)+(0,-1.35)$) node[below] {$\hat{z}_b$};
  \draw[->] (Ct)--($(Ct)+(-2.05,-2.00)$) node[below] {$\hat{y}_b$};
  \draw[->] (A)--($(A)+(-1.10,-0.55)$) node[left] {$\hat{x}_w$};
  \node at (6.05,2.65) {$+\beta$};
  \node at (6.45,1.85) {$+\alpha$};
\end{tikzpicture}
}
  \caption{Angle-of-attack and sideslip definitions for backward flight.}
  \label{fig:alpha-beta-backward}
\end{figure}

When $\hat x_w\mathbin{\cdot}\hat x_b=0$, the wind vector lies in the
$\hat y_b\hat z_b$ plane. Equations~\eqref{eq:projected-angle-of-attack} and
\eqref{eq:projected-sideslip} then give $\alpha=\pm90^\circ$ and
$\beta=\pm90^\circ$ for every vector in that plane and therefore do not uniquely
define its direction. TAOS uses the special convention illustrated in
\cref{fig:aero-angles-at-90-alpha}, retaining
\cref{eq:projected-angle-of-attack} for $\alpha$ and replacing the sideslip
relation by
\begin{equation}
  \tan\beta
  =\frac{\hat x_w\mathbin{\cdot}\hat y_b}
  {\hat x_w\mathbin{\cdot}\hat z_b}.
  \label{eq:projected-sideslip-at-90-alpha}
\end{equation}

\begin{figure}[htbp]
  \centering
  \resizebox{0.66\linewidth}{!}{\begin{tikzpicture}[scale=0.95,>=Latex,line cap=round,line join=round]
  \coordinate (O) at (0,0);
  \draw[->] (O)--(-2.50,0) node[left] {$\hat{y}_b$};
  \draw[->] (O)--(0,-2.50) node[below] {$\hat{z}_b$};
  \draw[->,line width=1pt] (O)--(-1.65,-1.75) node[below left] {$\hat{x}_w$};
  \draw[->,line width=1pt] (O)--(1.65,-1.75) node[below right] {$\hat{x}_w$};
  \draw[->] (-0.95,-1.35) arc[start angle=-130,end angle=-92,radius=0.75];
  \draw[->] (0.95,-1.35) arc[start angle=-50,end angle=-88,radius=0.75];
  \node at (-0.85,-1.80) {$+\beta$};
  \node at (0.85,-1.80) {$-\beta$};
  \draw[->] (2.2,-0.02)--(2.2,1.20);
  \draw[->] (2.2,-0.02)--(2.2,-1.20);
  \node[right] at (2.25,1.15) {$\alpha=-90^{\circ}$};
  \node[right] at (2.25,0) {$\alpha=0^{\circ}$};
  \node[right] at (2.25,-1.15) {$\alpha=+90^{\circ}$};
  \node[align=left] at (-1.95,1.30) {Angle of attack and sideslip when\\$\hat{x}_w\cdot\hat{x}_b=0$};
\end{tikzpicture}
}
  \caption{Aerodynamic angles at a $90^\circ$ angle of attack.}
  \label{fig:aero-angles-at-90-alpha}
\end{figure}

Because the basic projection definitions are singular in this case, angle of attack
with sideslip is not recommended for trajectories whose wind vector lies in the
$\hat y_b\hat z_b$ plane. The pair $\alpha$ and $\beta_E$, or the pair
$\alpha_T$ and $\phi_w$, is preferable.

\taossubhead{Total angle of attack and windward meridian}

The total angle of attack $\alpha_T$ is the angle between $\hat x_b$ and
$\hat x_w$. The windward meridian $\phi_w$ is the angle between $\hat z_b$ and
the projection of $\vec V_{w\oplus}$ onto the $\hat y_b\hat z_b$ plane, as shown
in \cref{fig:total-alpha-windward-meridian}. Their definitions are
\begin{equation}
  \tan\alpha_T
  =\frac{\sqrt{(\hat x_w\mathbin{\cdot}\hat y_b)^2
                  +(\hat x_w\mathbin{\cdot}\hat z_b)^2}}
  {\hat x_w\mathbin{\cdot}\hat x_b},
  \label{eq:total-angle-of-attack}
\end{equation}
\begin{equation}
  \tan\phi_w
  =-\frac{\hat x_w\mathbin{\cdot}\hat y_b}
  {\hat x_w\mathbin{\cdot}\hat z_b}.
  \label{eq:windward-meridian}
\end{equation}

\begin{figure}[htbp]
  \centering
  \resizebox{0.78\linewidth}{!}{\input{manual/figures/chapter02/fig-2-15}}
  \caption{Total angle of attack and windward meridian definitions.}
  \label{fig:total-alpha-windward-meridian}
\end{figure}

The total angle of attack is nonnegative and ranges from $0^\circ$ to $180^\circ$.
The windward meridian may take any value but is normally expressed between
$0^\circ$ and $360^\circ$.

\taossubhead{Aerodynamic-angle relationships}

TAOS guidance rules specify one pair of aerodynamic angles and derive the remaining
angles from that pair. If $\alpha$ and $\beta_E$ are known, then
\begin{equation}
  \tan\alpha_T
  =\frac{\sqrt{\sin^2\beta_E
      +\cos^2\beta_E\sin^2\alpha}}
  {\cos\alpha\cos\beta_E},
  \label{eq:alpha-betae-to-total-alpha}
\end{equation}
\begin{equation}
  \tan\phi_w
  =-\frac{\sin\beta_E}
  {\sin\alpha\cos\beta_E},
  \label{eq:alpha-betae-to-windward-meridian}
\end{equation}
\begin{equation}
  \tan\beta
  =\frac{\sin\beta_E}
  {\cos\alpha\cos\beta_E}.
  \label{eq:alpha-betae-to-beta}
\end{equation}
For the special case $\alpha=\pm90^\circ$, TAOS uses $\beta=-\phi_w$.

If $\alpha$ and $\beta$ are known, then
\begin{equation}
  \tan\alpha_T=\sqrt{\tan^2\alpha+\tan^2\beta},
  \label{eq:alpha-beta-to-total-alpha}
\end{equation}
\begin{equation}
  \tan\phi_w=-\frac{\tan\beta}{\tan\alpha},
  \label{eq:alpha-beta-to-windward-meridian}
\end{equation}
\begin{equation}
  \sin\beta_E=-\sin\alpha_T\sin\phi_w.
  \label{eq:alpha-beta-to-euler-sideslip}
\end{equation}
For the special case $\alpha=\pm90^\circ$, TAOS uses $\phi_w=-\beta$.

Finally, if $\alpha_T$ and $\phi_w$ are known, then
\begin{equation}
  \tan\alpha
  =\frac{\sin\alpha_T\cos\phi_w}{\cos\alpha_T},
  \label{eq:total-alpha-phi-to-alpha}
\end{equation}
\begin{equation}
  \tan\beta
  =-\frac{\sin\alpha_T\sin\phi_w}{\cos\alpha_T},
  \label{eq:total-alpha-phi-to-beta}
\end{equation}
\begin{equation}
  \sin\beta_E=-\sin\alpha_T\sin\phi_w.
  \label{eq:total-alpha-phi-to-euler-sideslip}
\end{equation}
For the special case $\alpha_T=90^\circ$, TAOS uses $\beta=-\phi_w$.

After all aerodynamic angles have been determined,
\cref{eq:body-x-from-aero-angles,eq:body-y-from-aero-angles,eq:body-z-from-aero-angles}
are used to compute the body unit vectors. These calculations are performed by
\texttt{body\_attitude} in the \texttt{derivs} module. Conversely, when the body
unit vectors were obtained from Euler angles, \cref{eq:total-angle-of-attack,eq:windward-meridian} provide $\alpha_T$ and $\phi_w$, after which
\cref{eq:total-alpha-phi-to-alpha,eq:total-alpha-phi-to-beta,eq:total-alpha-phi-to-euler-sideslip} supply the other aerodynamic angles.

\taossubhead{Aerodynamic and bank-angle output variables}

\begin{longtable}{@{}>{\ttfamily}p{0.18\textwidth}p{0.18\textwidth}p{0.56\textwidth}@{}}
  alpha  & $\alpha$ & Angle of attack.\\
  alphat & $\alpha_T$ & Total angle of attack.\\
  bankgc & $\mu_{gc}$ & Geocentric bank angle.\\
  bankgd & $\mu_{gd}$ & Geodetic bank angle.\\
  beta   & $\beta$ & Sideslip angle.\\
  betae  & $\beta_E$ & Euler sideslip angle.\\
  phi    & $\phi_w$ & Windward meridian angle.\\
\end{longtable}
